\section{The transfer to the mirror base}\label{sec:legendre}

The relative tropical $2$-cycles of Theorem~\ref{thm:real} live on the base $B'_\sigma$, whereas the chains of
Theorem~\ref{thm:lift} are built over the base $B_Y$ of the degeneration of the mirror family. This section constructs a
piecewise linear homeomorphism $\mathfrak h\colon B'_\sigma\to B_Y$ that carries the discriminant of $B'_\sigma$ onto that of
$B_Y$ and the local system of integral tangent vectors of $B'_\sigma$ onto the dual of that of $B_Y$, and uses it to
transfer relative tropical $2$-cycles to $B_Y$ with the same boundary networks and the same node coefficients
(Proposition~\ref{prop:transfer}). The homeomorphism is assembled from three comparisons, each compatible with the
dual-pair complex of Section~\ref{sec:dualpair}: a change of fan-side subdivision (Proposition~\ref{prop:F}), by which $\Sigma'$ is
replaced by $\mathcal F_C$; a change of height (Lemma~\ref{lem:height}); and the dual-pair homeomorphism between the two bases built
from the same data (Proposition~\ref{prop:L1}). Lemmas~\ref{lem:D} and~\ref{lem:arc} describe the cell structure near a
node and near an arc of the discriminant; these are the neighbourhoods in which the local models of Section~\ref{sec:lift}
are placed.

\subsection{Change of fan-side subdivision and of height}\label{ssec:leg-F}
Three comparisons of bases occur in this paper: a change of fan-side subdivision (Proposition~\ref{prop:F}), a change of height
(Lemma~\ref{lem:height}) and the passage to the mirror base (Section~\ref{ssec:leg-transfer}). Each consists of a
piecewise linear homeomorphism of pairs and an isomorphism of local systems, and each carries relative tropical
$2$-cycles by the following lemma.

\begin{lemma}[transport of relative cycles]\label{lem:transport}
For $i=1,2$ let $B_i$ be a closed piecewise linear three-manifold, $\Gamma_i\subset B_i$ a compact one-dimensional
subpolyhedron, $\Lambda_i$ a local system of free abelian groups of rank three on $B_i\setminus\Gamma_i$, and suppose that
$\Gamma_i$ is the union of finitely many points and open arcs $A$, each carrying a flat section $d_A$ of $\Lambda_i$ over
$V_A\setminus\Gamma_i$ for a neighbourhood $V_A$ of the closure $\bar A$. For triangulations
$\mathfrak K$ of $B_i$ in which $\Gamma_i$ is a full subcomplex define $F_i$, $\mathcal D_i(E)=\Z d_A$ for $E\subseteq A$,
and $Z_i(\mathfrak K)$ as in Definition~\ref{def:relcycle}. Let $f\colon B_1\to B_2$ be a piecewise linear homeomorphism
with $f(\Gamma_1)=\Gamma_2$, mapping each arc $A$ onto an arc $f(A)$, and $\theta\colon\Lambda_1\to f^*\Lambda_2$ an
isomorphism of local systems on $B_1\setminus\Gamma_1$ with $\theta(d_A)=d_{f(A)}$ for every arc $A$. Let $\mathfrak K$ be a
triangulation of $B_1$ in which $\Gamma_1$ is full, $\mathfrak K^0_1$ a subdivision of $\mathfrak K$ on whose simplices $f$
is linear, and $\mathfrak K_1$ a derived subdivision of $\mathfrak K^0_1$. Then:
\begin{enumerate}
\item $f(\mathfrak K_1)$ is a triangulation of $B_2$ in which $\Gamma_2$ is a full subcomplex; if $f$ is linear on the
simplices of $\mathfrak K$, the same holds for $f(\mathfrak K)$, and one may take $\mathfrak K_1=\mathfrak K$ below;
\item $\theta$ induces isomorphisms $F_1(s)\to F_2(f(s))$, compatible with the restrictions and carrying $\mathcal
D_1(E)$ onto $\mathcal D_2(f(E))$, hence an isomorphism of chain complexes $f_\#\colon C_*(\mathfrak K_1;F_1)\to
C_*(f(\mathfrak K_1);F_2)$ mapping $C_*(\Gamma_1;\mathcal D_1)$ onto $C_*(\Gamma_2;\mathcal D_2)$;
\item the composite $f_*:=f_\#\circ\lgsd\colon Z_1(\mathfrak K)\to Z_2(f(\mathfrak K_1))$, with $\lgsd$ the subdivision
operator, is an injective homomorphism, and if $\partial z=\sum_Ec_Ed_{A(E)}[E]$, $A(E)$ the arc containing the edge $E$,
then $\partial f_*z=\sum_Ec_E\sum_{E'\subseteq E}d_{f(A(E))}[f(E')]$, the inner sum over the edges $E'$ of $\mathfrak K_1$
in $E$ with the orientation of $E$.
\end{enumerate}
\end{lemma}

\begin{proof}
(1) The subcomplex $\Gamma_1$ of $\mathfrak K$ is a subcomplex of every subdivision, and in a derived subdivision every
subcomplex is full \cite[Ch.~3]{RS}. As $f$ is linear on the simplices of $\mathfrak K_1$, it is a simplicial isomorphism
onto $f(\mathfrak K_1)$, which is therefore a triangulation of $B_2$ with $\Gamma_2=f(\Gamma_1)$ full. The derived
subdivision is needed: in an arbitrary subdivision on which $f$ is linear, a simplex with all its vertices on $\Gamma_1$
need not lie in $\Gamma_1$.

(2) For a simplex $s$ of $\mathfrak K_1$, $f$ maps the neighbourhoods $V$ of $\bar s$ onto neighbourhoods $f(V)$ of
$f(\bar s)$, all small ones arising in this way, with $f(V\setminus\Gamma_1)=f(V)\setminus\Gamma_2$, and $\theta$ maps the
flat sections of $\Lambda_1$ over $V\setminus\Gamma_1$ to the flat sections of $\Lambda_2$ over $f(V)\setminus\Gamma_2$,
compatibly with restriction to smaller sets. By hypothesis it carries $d_A$ to $d_{f(A)}$.

(3) The subdivision operator sends an oriented simplex $s$ of $\mathfrak K$ with coefficient $X\in F(s)$ to the sum of the
simplices $s'$ of $\mathfrak K_1$ of the same dimension in $\bar s$, with the induced orientations and the coefficients
$X|_{s'}$; the restriction is defined because a neighbourhood of $\bar s$ is a neighbourhood of $\bar s'$. It is a chain
map, it maps $C_*(\Gamma_1;\mathcal D_1)$ into itself, since the edges of $\mathfrak K_1$ in an edge $E\subseteq A$ of
$\Gamma_1$ lie in $A$, and it is injective, since restriction of flat sections to a smaller connected set is injective and
$V\setminus\Gamma_1$ is connected for small $V$. So it maps $Z_1(\mathfrak K)$ injectively into $Z_1(\mathfrak K_1)$, with
the stated boundary, and (2) applies.
\end{proof}

\begin{proposition}[change of fan-side subdivision]\label{prop:F}
Let $\mathcal F_1$ and $\mathcal F_2$ be two fan-side subdivisions which agree on the two-skeleton of $\partial\Delta$: a cell of either
contained in a face of $\Delta$ of dimension at most two is a cell of the other. Write $C^i$, $\Gamma^i_X$, $\Lambda^i_X$ for
the cells, discriminant and local system of $B_X(\mathcal F_i)$; the underlying space $\partial\nabla^{\check h}$ is the
same.
\begin{enumerate}
\item The dual pairs $(\tau,\omega)$ with $\dim\tau\ge1$ are the same for $\mathcal F_1$ and $\mathcal F_2$.
\item If $\omega$ is not contained in a face of $\Delta$ of dimension at most two, then every dual pair $(\tau,\omega)$ has
$\tau=\{u_G\}$, where $G$ is the facet of $\Delta$ containing $\omega$ and $u_G$ the vertex of $\Delta^\circ$ dual to $G$.
\item There is a piecewise linear homeomorphism $g$ of $\partial\nabla^{\check h}$, preserving every face $F_\tau$ and
isotopic to the identity through face-preserving homeomorphisms, such that $g(C^1(\tau,\omega))=C^2(\tau,\omega)$
for $\dim\tau\ge1$. In particular $g(\Gamma^1_X)=\Gamma^2_X$, $g$ maps the arc $C^1([u,u'],[n,n'])$ onto the arc
$C^2([u,u'],[n,n'])$ and each vertex of $\Gamma^1_X$ onto the vertex of $\Gamma^2_X$ of the same dual pair, and there is an
isomorphism $\kappa\colon\Lambda^1_X\to g^*\Lambda^2_X$ that is the identity of $u^\perp\cap N$ on every open facet
$\operatorname{int}F_u$ and carries the vanishing vector $n'-n$ of each arc to that of its image.
\item If $h_1$ and $h_2$ agree on the lattice points of the two-skeleton, $g$ can be taken to be the identity on the
two-skeleton of $\partial\nabla^{\check h}$.
\end{enumerate}
\end{proposition}
\begin{proof}
(1) If $\dim\tau\ge1$, then $\omega\subseteq\tau^\vee$, a face of $\Delta$ of dimension at most two, where the two
subdivisions agree. (2) The smallest face of $\Delta$ containing $\omega$ is a facet $G$, and $\tau^\vee$ is a proper face
of $\Delta$ containing $\omega$, so $\tau^\vee=G$; the only point $m\in\Delta^\circ$ with $\langle m,\cdot\,\rangle=-1$ on
$G$ is $u_G$.

(3) By Proposition~\ref{prop:L1}(5), the two-skeleton $\bigcup_{\dim\tau\ge1}F_\tau$ is the union of the cells
$C^i(\tau,\omega)$ with $\dim\tau\ge1$; these form a subcomplex of each dual-pair complex, with the same indexing set and
order by (1). The inductive construction in the proof of Proposition~\ref{prop:L1}(4), applied to these two regular CW
structures on the same space, gives a piecewise linear homeomorphism $g$ of the two-skeleton with
$g(C^1(\tau,\omega))=C^2(\tau,\omega)$; it preserves every $F_\tau$ with $\dim\tau\ge1$, which is the union of the cells
$C^i(\tau',\omega)$ with $\tau'\supseteq\tau$. Extend $g$ over each facet, a convex three-ball whose boundary it preserves,
by coning from an interior point; the cone of a piecewise linear map is piecewise linear. Faces of dimension zero are
points and are fixed; on the edges, the two-faces and the facets in turn, $g$ is isotopic to the identity relative to the
lower-dimensional faces by the Alexander trick, extended by coning as in Appendix~\ref{app:leg-pl}. By
Propositions~\ref{prop:L2}(1),(2), $g(\Gamma^1_X)=\Gamma^2_X$, with arcs and vertices as stated. Define $\kappa$ on each
$\operatorname{int}F_u$, which $g$ preserves, as the identity of $u^\perp\cap N$. A point $b$ of the two-skeleton outside
$\Gamma^1_X$ lies in some $C^1(\tau,n)$ with $\dim\tau\ge1$, so $b\in\widetilde R^1_n$ and $g(b)\in C^2(\tau,n)\subseteq
\widetilde R^2_n$. On a connected neighbourhood $U\subseteq\widetilde R^1_n\cap g^{-1}(\widetilde R^2_n)$ of $b$ both local
systems are the constant $N/\Z n$, and on $U\cap\operatorname{int}F_u$ the identity of $u^\perp\cap N$ corresponds to the
identity of $N/\Z n$ under $\pi^X_{un}$ on both sides. So $\kappa$ extends over $U$ as the identity of $N/\Z n$, compatibly
on overlaps, and is an isomorphism of local systems on $\partial\nabla^{\check h}\setminus\Gamma^1_X$. Near the arc
$C^i([u,u'],[n,n'])$ the vanishing vector is represented in the chart of $F_u$ by $n'-n\in u^\perp\cap N$ on both bases
(Proposition~\ref{prop:L2}(4)), and $\kappa$ is the identity there.

(4) By Lemma~\ref{lem:cayley}(2), the decomposition of a face $F_\tau$ with $\dim\tau\ge1$ depends only on the restriction
of $h_i$ to $\tau^\vee\cap N$, which lies in the two-skeleton of $\partial\Delta$. So the two decompositions, their
barycentric subdivisions and their cells agree on the two-skeleton, and $g$ can be taken to be the identity there.
\end{proof}

The vertices of $\mathcal F_2$ in the interior of a facet $G$ of $\Delta$ only change the cells $C^2(u_G,\omega)$ inside the
open facet $\operatorname{int}F_{u_G}$, by (2); there $\Lambda^2_X$ is the constant facet lattice, and $\kappa$ is its
identity.

We apply Proposition~\ref{prop:F} with $\mathcal F_1=\Sigma'$ and $\mathcal F_2=\mathcal F_C$. The pulling refinement
$\mathcal F_C$ is a fan-side subdivision, induced by the height $h_C$, and it agrees with $\Sigma'$ on the two-skeleton
(Proposition~\ref{prop:pulling}). We build $B_X(\mathcal F_C)$ with the function $h_C$, and fix the resulting piecewise
linear homeomorphism $g_C\colon B'_\sigma\to B_X(\mathcal F_C)$ and the isomorphism $\kappa_C\colon\Lambda\to g_C^*\Lambda^{\mathcal F_C}_X$. Since $g_C$ maps the vertex $p_q$ of the discriminant,
the zero-cell $C([u,u'],P_q)$, to the zero-cell of the same dual pair, and each arc to the arc with the same lower support,
the node coefficients of Definition~\ref{def:nodecoef}, read from the lower supports of the legs at $p_q$, are preserved by
the transport $(g_C)_*$ of Lemma~\ref{lem:transport}.

\begin{remark}\label{rem:gCid}
If every lattice point of $\partial\Delta$ is a vertex of $\Sigma'$, no point is pulled and $\mathcal F_C=\Sigma'$; if in
addition $B_X(\mathcal F_C)$ is built with the same function $h$ as $B'_\sigma$, then $g_C$ can be taken to be the identity.
This is the case for the subdivisions of the examples of Section~\ref{sec:examples}. Nothing in this paper uses it.
\end{remark}

\begin{lemma}[change of height]\label{lem:height}
Let $\check h_1$ and $\check h_2$ be two integral heights on the same unimodular triangulation $\mathcal T$ of
$\partial\Delta^\circ$, both with $\check h_i$ and $\check h_i-\check\varphi$ strictly convex on the fan over $\mathcal T$, and
let $B_i=\partial\nabla^{\check h_i}$ be the base $B_X(\Sigma')$ built from $\check h_i$, with the same function $h$ in
$\widetilde\Sigma'$. Write $C^i(\tau,\omega)$ for the cells of its dual-pair complex, $\Gamma_i$ for its discriminant and
$\Lambda_i$ for its local system. Then:
\begin{enumerate}
\item there is a piecewise linear homeomorphism $\phi\colon B_1\to B_2$ with $\phi(C^1(\tau,\omega))=C^2(\tau,\omega)$ for
every dual pair $(\tau,\omega)$; it maps each open facet $\operatorname{int}F_u$ and each region $\widetilde R_n$ onto the
open facet and the region with the same index, $\Gamma_1$ onto $\Gamma_2$, each vertex $p_q$ of the discriminant onto $p_q$, each arc onto
the arc with the same dual pair, and each vertex of $\Gamma_1$ onto a vertex of $\Gamma_2$ of the same type;
\item there is an isomorphism of local systems $\kappa\colon\Lambda_1\to\phi^*\Lambda_2$ on $B_1\setminus\Gamma_1$ that is the
identity of $u^\perp\cap N$ on every $\operatorname{int}F_u$ and the identity of $N/\Z n$ on every $\widetilde R_n$;
\item the transport $\phi_*$ of Lemma~\ref{lem:transport} carries relative tropical $2$-cycles on $B_1$ to relative
tropical $2$-cycles on $B_2$ with the same integers $c_E$ in their boundary networks and the same node coefficients; and
$(S,j,v)\mapsto(S,\phi\circ j,\kappa(v))$ is a bijection from the tropical $2$-cycles of \cite[Def.~7.2]{CBM} on $B_1$ to
those on $B_2$, which preserves piecewise linearity of $j$ and the node coefficients.
\end{enumerate}
In particular, relative tropical $2$-cycles and tropical $2$-cycles on the base of any such height on $\mathcal T$ are
carried to cycles of the same kind on the base of the normalised height of Section~\ref{ssec:gross}, with the same node
coefficients.
\end{lemma}
\begin{proof}
(1) The homeomorphism of Lemma~\ref{lem:indepheight} has these properties. Alternatively: the set of dual pairs and its order depend only on $\mathcal T$ and $\Sigma'$, and Proposition~\ref{prop:L1}, which
uses only the strict convexity of $\check h_i$ and $\check h_i-\check\varphi$, applies to both heights. The inductive
construction in its proof gives piecewise linear homeomorphisms $\Psi_i\colon|\lgOrd(\lgDP)|\to B_i$ with
$\Psi_i(|\lgOrd^{\preceq {\mathfrak p}}|)=\overline{C^i({\mathfrak p})}$, and $\phi=\Psi_2\circ\Psi_1^{-1}$ preserves the cells. The facets, regions,
arcs and vertices of the discriminant are unions of cells described by their dual pairs (Propositions~\ref{prop:L1}(5)
and~\ref{prop:L2}(1),(2)); the vertex $p_q$ of the discriminant is the zero-cell $C^i([u,u'],P_q)$, the barycentre of the carrying cell, the
unique cell with face label $[u,u']$ and lower support $P_q$ (Lemma~\ref{lem:stack}); and the type of a three-valent vertex
is determined by its dual pair (Proposition~\ref{prop:L2}(3)).

(2) The open facets and the regions cover $B_i\setminus\Gamma_i$, with $\operatorname{int}F_u\cap\widetilde R_n=C^i(u,n)$
and no triple intersections (proof of Proposition~\ref{prop:L3}). On $\operatorname{int}F_u$ the lattice $\Lambda_i$ is
$u^\perp\cap N$, on $\widetilde R_n$ it is $N/\Z n$, and on $C^i(u,n)$ the two are related by the projection
$u^\perp\cap N\to N/\Z n$, on both bases. So the identities on the pieces agree on the overlaps and define $\kappa$.

(3) The pair $(\phi,\kappa)$ satisfies the hypotheses of Lemma~\ref{lem:transport}: near the arc of $([u,u'],[n,n'])$ the
vanishing vector is $n'-n\in u^\perp\cap N$ on both bases and $\kappa$ is the identity there. The integers $c_E$ are
therefore unchanged, and the node coefficients, read from the lower supports of the legs at the node points
(Definition~\ref{def:nodecoef}), are unchanged by (1). The clauses of \cite[Def.~7.2]{CBM} refer to the topology of $S$
and $j$, to $\Gamma_i$ with its vertex types and nodes, and to $\Lambda_i$ and its monodromy. An orientation of $B_i$ enters
only through Definition~\ref{def:types}, and the types do not depend on it (Remark~\ref{rem:sameloop}). By (1) and (2),
$\phi$ and $\kappa$ preserve all of these, and $\kappa$ intertwines parallel transport, so $(S,\phi\circ j,\kappa(v))$
satisfies the clauses if and only if $(S,j,v)$ does; $\phi\circ j$ is piecewise linear if $j$ is. For a tropical
$2$-cycle, triangulated, the node coefficients are those of the relative cycle it defines (Section~\ref{ssec:cbm}), and
these are unchanged.
\end{proof}

\subsection{The node star and the neighbourhood of an arc}\label{ssec:leg-node}
In this subsection $B_X=B_X(\mathcal F_C)$, $B_Y$ is the mirror base, and $\mathcal L=\Psi_Y\circ\Psi_X^{-1}$ is the
dual-pair homeomorphism constructed in the proof of Proposition~\ref{prop:L1}(4), which we fix. Since $g_C$ preserves the faces of
$\partial\nabla^{\check h}$ and the cells $C(\tau,\omega)$ with $\dim\tau\ge1$ (Proposition~\ref{prop:F}(3)), statements
about $B'_\sigma$ that involve only these are carried over by $g_C$.

Let $q\in Q$ be a node, with node parallelogram $P=abcd$ in the two-face $G$ of $\Delta$, labelled as in
Section~\ref{ssec:admissible}, and $G^\circ=[u,u']$, an edge of $\mathcal T$ because $\ell(G^\circ)=1$. Then
$([u,u'],P)$ is a dual pair of dimension zero; its zero-cell in $B'_\sigma$ is the vertex $p_q$ of the discriminant, the barycentre of the
carrying cell, the unique cell of $\mathscr P$ with face label $[u,u']$ and lower support $P$ (Lemma~\ref{lem:stack}).
Write ${\mathfrak q}=([u,u'],P)$ for this dual pair; its zero-cells are indexed by the node $q$. Put $f_1=b-a=c-d$ and $f_2=d-a=c-b$. The statements
below concern only the cells of $\mathcal T$ contained in $[u,u']$ and the cells of $\mathcal F_C$ contained in $P$; they do
not depend on whether $G$ is a square, a pentagon or a hexagon, nor, for a hexagon, on the branch of the profile.

To a face $\omega$ of $P$ assign $(\epsilon_1(\omega),\epsilon_2(\omega))\in\{-,0,+\}^2$, its position in the directions
$f_1$ and $f_2$: $a\mapsto(-,-)$, $b\mapsto(+,-)$, $c\mapsto(+,+)$, $d\mapsto(-,+)$, $ab\mapsto(0,-)$, $bc\mapsto(+,0)$,
$cd\mapsto(0,+)$, $da\mapsto(-,0)$, $P\mapsto(0,0)$. To $\tau\in\{u,[u,u'],u'\}$ assign $\epsilon_3(u)=-$,
$\epsilon_3([u,u'])=0$, $\epsilon_3(u')=+$. In $\R^3$ with coordinates $(y_1,y_2,y_3)$ let
$\mathcal O_\epsilon=\{y:\operatorname{sign}y_i=\epsilon_i\}$ for $\epsilon\in\{-,0,+\}^3$, the relatively open coordinate
orthants of all dimensions, and $O=\{|y_1|+|y_2|+|y_3|\le1\}$. Put $\operatorname{St}^X_q=\Psi_X(\operatorname{st}{\mathfrak q})$ and
$\operatorname{St}^Y_q=\Psi_Y(\operatorname{st}{\mathfrak q})$, where $\operatorname{st}{\mathfrak q}$ is the closed star of the vertex ${\mathfrak q}$ in
$\lgOrd(\lgDP)$; these are closed neighbourhoods of the zero-cells indexed by $q$, and $\mathcal L(\operatorname{St}^X_q)=\operatorname{St}^Y_q$.
Finally let $\mathcal W_q$ be the union of the cells $C_Y([u,u'],\omega)$, $\omega\subseteq P$: the node, the four arcs at it
and the four two-cells between them.

\begin{lemma}[the node star]\label{lem:D}
\begin{enumerate}
\item The dual pairs ${\mathfrak p}\succeq{\mathfrak q}$ are the $27$ pairs $(\tau,\omega)$ with $\tau$ equal to $u$, $[u,u']$ or $u'$ and
$\omega$ a face of $P$: the node; four arcs, $\tau=[u,u']$ and $\omega$ an edge; four two-cells, $\tau=[u,u']$ and
$\omega$ a vertex; two one-cells, $\omega=P$ and $\tau$ a vertex; eight two-cells, $\tau$ a vertex and $\omega$ an edge;
and eight three-cells, $\tau$ and $\omega$ vertices. The map
$\epsilon(\tau,\omega)=(\epsilon_1(\omega),\epsilon_2(\omega),\epsilon_3(\tau))$ is an isomorphism of this poset onto the
face poset of the coordinate fan: ${\mathfrak p}\preceq {\mathfrak p}'$ if and only if $\mathcal O_{\epsilon({\mathfrak p})}\subseteq
\overline{\mathcal O_{\epsilon({\mathfrak p}')}}$; and $\dim{\mathfrak p}$ is the number of nonzero entries of $\epsilon({\mathfrak p})$.
\item There are piecewise linear homeomorphisms $\Phi^X_q\colon\operatorname{St}^X_q\to O$ and
$\Phi^Y_q\colon\operatorname{St}^Y_q\to O$ with $\Phi^X_q(C_X({\mathfrak p})\cap\operatorname{St}^X_q)=\mathcal O_{\epsilon({\mathfrak p})}\cap O=
\Phi^Y_q(C_Y({\mathfrak p})\cap\operatorname{St}^Y_q)$ for all ${\mathfrak p}\succeq{\mathfrak q}$, and $\Phi^Y_q\circ\mathcal L=\Phi^X_q$ on
$\operatorname{St}^X_q$. Any two homeomorphisms $\operatorname{St}^Y_q\to O$ with this property are isotopic through such
homeomorphisms.
\item Consequently $\Phi^Y_q(\mathcal W_q\cap\operatorname{St}^Y_q)=\{y_3=0\}\cap O$; $\Phi^Y_q$ maps
$\widetilde R^*_u\cap\operatorname{St}^Y_q$ and $\widetilde R^*_{u'}\cap\operatorname{St}^Y_q$ onto $\{y_3<0\}\cap O$ and
$\{y_3>0\}\cap O$, the open facets $\operatorname{int}F^*_a,\dots,\operatorname{int}F^*_d$ onto the four quadrants
$\{(\operatorname{sign}y_1,\operatorname{sign}y_2)=(\epsilon_1,\epsilon_2)\}$ of the vertices, and the edge $F^*_P$ of
$\nabla^H$ onto the $y_3$-axis. On $B_X$, $\Phi^X_q$ maps $F_{uu'}\cap\operatorname{St}^X_q$ onto $\{y_3=0\}\cap O$ and
$\operatorname{int}F_u$, $\operatorname{int}F_{u'}$ onto the two half-balls.
\end{enumerate}
\end{lemma}
\begin{proof}
(1) The pairs $\succeq{\mathfrak q}$ are the $(\tau,\omega)$ with $\tau\subseteq[u,u']$ and $\omega\subseteq P$, and every such pair is
dual, being contained in a dual pair. The faces of the parallelogram $P$ are the products of faces of two segments, and the
faces of the segment $[u,u']$ are its end points and itself; in each factor, inclusion of faces corresponds to
$\epsilon_i\in\{0,\epsilon'_i\}$ with $0$ for the larger face, which is the order of the coordinate fan. The number of
nonzero entries is $3-\dim\tau-\dim\omega$.

(2) The closed star of ${\mathfrak q}$ in $\lgOrd(\lgDP)$ consists of the chains containing ${\mathfrak q}$ and their subchains; since ${\mathfrak q}$ is
minimal in $\{{\mathfrak p}\succeq{\mathfrak q}\}$ and has $\dim{\mathfrak q}=0$, it is the cone ${\mathfrak q}*\lgOrd(\lgDP_{\succ{\mathfrak q}})$. The nonzero sign vectors, ordered
as in (1), form the face poset of the boundary of the cross-polytope $O$, the face of $\epsilon$ being
$\operatorname{conv}\{\epsilon_ie_i:\epsilon_i\ne0\}$; so $\epsilon$ induces a simplicial isomorphism of
$\lgOrd(\lgDP_{\succ{\mathfrak q}})$ with the barycentric subdivision of $\partial O$, whose cone is a piecewise linear homeomorphism
$E_q\colon|\operatorname{st}{\mathfrak q}|\to O$. The open simplices of $\operatorname{st}{\mathfrak q}$ whose chain has top ${\mathfrak p}$ are mapped onto
$\mathcal O_{\epsilon({\mathfrak p})}\cap O$, and $\Psi_X$ maps their union onto $C_X({\mathfrak p})\cap\operatorname{St}^X_q$, because open
simplices of chains with top ${\mathfrak p}$ lie in $\Psi_X^{-1}(C_X({\mathfrak p}))$. So $\Phi^X_q=E_q\circ\Psi_X^{-1}$ and
$\Phi^Y_q=E_q\circ\Psi_Y^{-1}$ have the stated properties, and $\Phi^Y_q\circ\mathcal L=\Phi^X_q$. Two homeomorphisms with
the property differ by a homeomorphism of $O$ preserving every $\mathcal O_\epsilon\cap O$, which is isotopic to the
identity through such homeomorphisms (Appendix~\ref{app:leg-pl}).

(3) The cells of $\mathcal W_q$ are those with $\tau=[u,u']$, that is, $\epsilon_3=0$. By
Proposition~\ref{prop:L1}(5), $\widetilde R^*_u\cap\operatorname{St}^Y_q$ is the union of the cells $C_Y(u,\omega)$, that
is, $\epsilon_3=-$; $\operatorname{int}F^*_a\cap\operatorname{St}^Y_q$ is the union of the cells $C_Y(\tau,a)$, that is,
$(\epsilon_1,\epsilon_2)=(-,-)$; and $F^*_P\cap\operatorname{St}^Y_q$ is the union of the cells $C_Y(\tau,P)$, that is,
$\epsilon_1=\epsilon_2=0$. On $B_X$, $F_{uu'}$ and $\operatorname{int}F_u$ are the unions of the cells with $\tau=[u,u']$,
respectively $\tau=u$, by Proposition~\ref{prop:L1}(5).
\end{proof}

The arcs at the node map under $\Phi^Y_q$ to the four half-axes in $\{y_3=0\}$: the arc over $ab$ to
$\{y_1=y_3=0,\ y_2<0\}$, over $dc$ to $y_2>0$, over $da$ to $y_1<0$ and over $bc$ to $y_1>0$. The four \emph{standard
halves} at $q$ are the sets $(\Phi^Y_q)^{-1}(\{y_1=0,\ \pm y_3\ge0\})$ and $(\Phi^Y_q)^{-1}(\{y_2=0,\ \pm y_3\ge0\})$; by (1)
and (3), $(\Phi^Y_q)^{-1}(\{y_1=0,\ y_3\ge0\})$ is the union of the closed cells with sign vectors $(0,\ast,+)$ and
$(0,\ast,0)$, that is, the closure of $C_Y(u',ab)\cup C_Y(u',P)\cup C_Y(u',cd)$, and similarly for the other three. In the
node model of Section~\ref{sec:regime}, the tropical chart $T_q$ of Lemma~\ref{lem:nodemodel}(4) has the same
quadrant and half-space description. That lemma compares the chart with scaled logarithmic coordinates on the
hypersurface up to $O(1/L)$; the cell identification here concerns the tropical bases.

Let ${\mathfrak p}_e=([u,u'],[n,n'])$ be a dual pair of edges, so that $C_Y({\mathfrak p}_e)$ is an arc of $\Gamma_Y$
(Proposition~\ref{prop:L2}(2)). To the faces of $[n,n']$ assign $\epsilon_1(n)=-$, $\epsilon_1(n')=+$,
$\epsilon_1([n,n'])=0$, and to the faces of $[u,u']$ assign $\epsilon_2(u)=-$, $\epsilon_2(u')=+$, $\epsilon_2([u,u'])=0$.
For $\epsilon\in\{-,0,+\}^2$ let $\mathcal O^{(2)}_\epsilon=\{(x_1,x_2):\operatorname{sign}x_i=\epsilon_i\}$. In $\R^3$ with
coordinates $(x_1,x_2,x_3)$ let $V=\operatorname{conv}\bigl((\Diamond\times\{0\})\cup\{\pm e_3\}\bigr)$, where
$\Diamond=\{|x_1|+|x_2|\le1\}$, and $V^\circ=V\setminus\{\pm e_3\}$. Put $\operatorname{St}^X_e=\Psi_X(\operatorname{st}{\mathfrak p}_e)$
and $\operatorname{St}^Y_e=\Psi_Y(\operatorname{st}{\mathfrak p}_e)$, with $\operatorname{st}{\mathfrak p}_e$ the closed star of the vertex ${\mathfrak p}_e$ in
$\lgOrd(\lgDP)$; then $\mathcal L(\operatorname{St}^X_e)=\operatorname{St}^Y_e$.

\begin{lemma}[the neighbourhood of an arc]\label{lem:arc}
\begin{enumerate}
\item The dual pairs ${\mathfrak p}\succeq {\mathfrak p}_e$ are the nine pairs $(\tau,\omega)$ with $\tau\subseteq[u,u']$ and
$\omega\subseteq[n,n']$. The map $\epsilon(\tau,\omega)=(\epsilon_1(\omega),\epsilon_2(\tau))$ is an isomorphism of this
poset onto the face poset of the coordinate fan of $\R^2$, and $\dim{\mathfrak p}-1$ is the number of nonzero entries of $\epsilon({\mathfrak p})$.
There are exactly two dual pairs $z_\pm\prec {\mathfrak p}_e$; they have dimension zero, and $C_Y(z_\pm)$ are the end points of the arc.
\item There are piecewise linear homeomorphisms $\Phi^X_e\colon\operatorname{St}^X_e\to V$ and
$\Phi^Y_e\colon\operatorname{St}^Y_e\to V$ with $\Phi^Y_e\circ\mathcal L=\Phi^X_e$, mapping $C_Y(z_\pm)$ to $\pm e_3$, the arc
$C_Y({\mathfrak p}_e)$ onto the open segment $\{x_1=x_2=0,\ |x_3|<1\}$, and $C_Y({\mathfrak p})\cap\operatorname{St}^Y_e$ onto
$\{x\in V^\circ:(x_1,x_2)\in\mathcal O^{(2)}_{\epsilon({\mathfrak p})}\}$ for every ${\mathfrak p}\succ {\mathfrak p}_e$; likewise on $B_X$. The set
$\operatorname{St}^Y_e$ is a neighbourhood of the open arc.
\item Consequently $\Phi^Y_e$ maps $\operatorname{int}F^*_n\cap\operatorname{St}^Y_e$ and
$\operatorname{int}F^*_{n'}\cap\operatorname{St}^Y_e$ onto $\{x_1<0\}\cap V^\circ$ and $\{x_1>0\}\cap V^\circ$, $\widetilde
R^*_u\cap\operatorname{St}^Y_e$ and $\widetilde R^*_{u'}\cap\operatorname{St}^Y_e$ onto $\{x_2<0\}\cap V^\circ$ and
$\{x_2>0\}\cap V^\circ$, the two-cells $C_Y(u,[n,n'])$ and $C_Y(u',[n,n'])$ onto $\{x_1=0,\ x_2<0\}\cap V^\circ$ and
$\{x_1=0,\ x_2>0\}\cap V^\circ$, and $C_Y([u,u'],n)$, $C_Y([u,u'],n')$ onto $\{x_2=0,\ x_1<0\}\cap V^\circ$ and
$\{x_2=0,\ x_1>0\}\cap V^\circ$.
\item Suppose that an end of the arc is a node $q$, so that $[n,n']$ is an edge of the node parallelogram $P$, labelled so
that $n'-n\in\{f_1,f_2\}$. For ${\mathfrak p}\succeq {\mathfrak p}_e$ the sign vector of ${\mathfrak p}$ in Lemma~\ref{lem:D} is obtained from
$\epsilon({\mathfrak p})=(\epsilon'_1,\epsilon'_2)$ by inserting the constant sign of the arc: it is $(\epsilon'_1,-,\epsilon'_2)$,
$(\epsilon'_1,+,\epsilon'_2)$, $(-,\epsilon'_1,\epsilon'_2)$, $(+,\epsilon'_1,\epsilon'_2)$ for the arcs over $ab$, $dc$,
$ad$, $bc$. In particular, on $\operatorname{St}^Y_e\cap\operatorname{St}^Y_q$ the standard halves containing the arc, on
the sides $y_3\ge0$ and $y_3\le0$, are the closures in $\operatorname{St}^Y_e\cap\operatorname{St}^Y_q$ of $C_Y(u',[n,n'])$
and $C_Y(u,[n,n'])$, which $\Phi^Y_e$ maps to $\{x_1=0,\ x_2\ge0\}$ and $\{x_1=0,\ x_2\le0\}$.
\end{enumerate}
\end{lemma}
\begin{proof}
(1) The pairs $\succeq {\mathfrak p}_e$ are the $(\tau,\omega)$ with $\tau\subseteq[u,u']$ and $\omega\subseteq[n,n']$, all dual, being
contained in a dual pair. The faces of a segment are its two end points and itself, and in each factor inclusion
corresponds to $\epsilon_i\in\{0,\epsilon'_i\}$ with $0$ for the larger face; $\dim{\mathfrak p}=3-\dim\tau-\dim\omega$ is one more than
the number of end points among $\tau,\omega$. By Proposition~\ref{prop:L1}(2), $\overline{C_Y({\mathfrak p}_e)}$ is a one-ball whose
boundary is the union of the nonempty cells $C_Y(z)$, $z\prec {\mathfrak p}_e$; these have $d_z=0$, so they are points, and there are
exactly two of them.

(2) This is a join construction like that of Lemma~\ref{lem:D}(2); the details are in
Appendix~\ref{app:leg-pl}.

(3) By Proposition~\ref{prop:L1}(5), $\operatorname{int}F^*_n=\bigcup_\tau C_Y(\tau,n)$ and $\widetilde
R^*_u=\bigcup_\omega C_Y(u,\omega)$; within $\operatorname{St}^Y_e$ these are the cells with $\epsilon_1=-$, respectively
$\epsilon_2=-$. The other statements are (2) for the individual cells.

(4) The sign vector of Lemma~\ref{lem:D} is $(\epsilon_1(\omega),\epsilon_2(\omega),\epsilon_3(\tau))$ for the faces
$\omega$ of $P$. For $\omega\subseteq[a,b]$ the second entry is $-$ and the first is $\epsilon_1(\omega)$ of the present
lemma, with $(n,n')=(a,b)$; the other three arcs are the same, with $(n,n')=(d,c)$, $(a,d)$, $(b,c)$. The standard half
$\{y_1=0,\ y_3\ge0\}$ meets $\operatorname{St}^Y_e\cap\operatorname{St}^Y_q$, near the arc over $ab$, in the closure of
$\mathcal O_{(0,-,+)}$, the cell of $(u',[a,b])$, and similarly in the other cases; now apply (3).
\end{proof}

Lemmas~\ref{lem:D} and~\ref{lem:arc} are statements about the cells $C_Y({\mathfrak p})$ and hold for their images under any
homeomorphism $\Upsilon$ of $B_Y$: with $\Phi^\Upsilon_q:=\Phi^Y_q\circ\Upsilon^{-1}$ and
$\Phi^\Upsilon_e:=\Phi^Y_e\circ\Upsilon^{-1}$ they hold for the cells $\Upsilon(C_Y({\mathfrak p}))$. In Section~\ref{sec:lift} they are
used in this form for the model complex.

\subsection{Transfer of relative tropical \texorpdfstring{$2$}{2}-cycles}\label{ssec:leg-transfer}
The lift of Section~\ref{sec:lift} is built over $B_Y$, after the discriminant of $B_Y$ has been moved by a piecewise linear
homeomorphism $\Upsilon^M$ of $B_Y$ to the positions over which the local models of the mirror sit
(Lemma~\ref{lem:modelcx}). We therefore define the transfer for an arbitrary piecewise linear homeomorphism $\Upsilon$ of
$B_Y$. Of Lemma~\ref{lem:modelcx} only the fact that $\Upsilon^M$ is such a homeomorphism is used here; the cells, the
discriminant and the local system of the model complex are then the transported ones of Definition~\ref{def:transfer}.

\begin{definition}[the transfer]\label{def:transfer}
Let $\Upsilon$ be a piecewise linear homeomorphism of $B_Y$. Put $C^\Upsilon({\mathfrak p}):=\Upsilon(C_Y({\mathfrak p}))$ for ${\mathfrak p}\in\lgDP$,
$\Gamma^\Upsilon:=\Upsilon(\Gamma_Y)$ and $\Lambda^\Upsilon:=\Upsilon_*\Lambda_Y$, and let
\[
\mathfrak h:=\Upsilon\circ\mathcal L\circ g_C\colon B'_\sigma\longrightarrow B_Y,
\]
with $g_C$ of Section~\ref{ssec:leg-F} and the dual-pair homeomorphism $\mathcal L\colon B_X(\mathcal F_C)\to B_Y$ fixed in
Section~\ref{ssec:leg-node}. Let
\[
\kappa_{\mathfrak h}\colon\Lambda\longrightarrow \mathfrak h^*(\Lambda^\Upsilon)^\vee\qquad\text{on }B'_\sigma\setminus\Gamma
\]
be the composite of $\kappa_C\colon\Lambda\to g_C^*\Lambda^{\mathcal F_C}_X$ (Proposition~\ref{prop:F}(3)), of the
isomorphism $\Lambda^{\mathcal F_C}_X\to(\mathcal L^*\Lambda_Y)^\vee$, $x\mapsto\langle\,\cdot\,,x\rangle_{\mathcal L}$, given by the flat
perfect pairing of Proposition~\ref{prop:L3}, and of the transport by $\Upsilon$. For an arc
$A=C^\Upsilon([u,u'],[n,n'])$ of $\Gamma^\Upsilon$ put $d^*_A:=\langle\,\cdot\,,n'-n\rangle$, a flat section of
$(\Lambda^\Upsilon)^\vee$ near $A$ (Proposition~\ref{prop:dict}(2)). For a triangulation $\mathfrak K^*$ of $B_Y$ in which
$\Gamma^\Upsilon$ is full, let $F^*(s)$ be the group of flat sections of $(\Lambda^\Upsilon)^\vee$ over
$V\setminus\Gamma^\Upsilon$, $V$ a small neighbourhood of $\bar s$, let $\mathcal D^*(E):=\Z d^*_A$ for an edge
$E\subseteq A$ of $\Gamma^\Upsilon$, and let $Z(\mathfrak K^*;F^*,\mathcal D^*)$ be the group of chains
$z^*\in C_2(\mathfrak K^*;F^*)$ with $\partial z^*\in C_1(\Gamma^\Upsilon;\mathcal D^*)$, as in
Definition~\ref{def:relcycle}.

For a triangulation $\mathfrak K$ of $B'_\sigma$ in which $\Gamma$ is full, choose a subdivision $\mathfrak K^0_1$ of
$\mathfrak K$ on whose simplices $\mathfrak h$ is linear and a derived subdivision $\mathfrak K_1$ of $\mathfrak K^0_1$. The
\emph{transfer} of $z\in Z(\mathfrak K)$ is $\mathfrak h_*z:=\mathfrak h_\#(\lgsd z)$, the transport of Lemma~\ref{lem:transport} for the pair
$(\mathfrak h,\kappa_{\mathfrak h})$.
\end{definition}

The group $F^*(s)$ is the analogue on $B_Y$ of the coefficient group $F(s)$ of Definition~\ref{def:relcycle}, for the
dual local system: the coefficients of a transferred cycle are integral cotangent vectors of $B_Y$, invariant under the
local monodromy. In Section~\ref{sec:lift}, $\Upsilon=\Upsilon^M$, and we write $C^M({\mathfrak p})$, $\Gamma^M$ and $\Lambda^M$ for
$C^\Upsilon({\mathfrak p})$, $\Gamma^\Upsilon$ and $\Lambda^\Upsilon$. For $x\in\Gamma^\Upsilon$ let
$\operatorname{Inv}^*(x)$ be the group of flat sections of $(\Lambda^\Upsilon)^\vee$ over $V\setminus\Gamma^\Upsilon$ for a
small ball $V$ about $x$, and for an arc $A$ let $\operatorname{Inv}^*(A)$ be its value at the points of $A$.

\begin{proposition}[the transferred coefficients]\label{prop:dict}
\begin{enumerate}
\item The homeomorphism $\mathfrak h$ maps every cell $C({\mathfrak p})$ of $B'_\sigma$ with ${\mathfrak p}=(\tau,\omega)$, $\dim\tau\ge1$, onto
$C^\Upsilon({\mathfrak p})$. In particular $\mathfrak h(\Gamma)=\Gamma^\Upsilon$; $\mathfrak h$ maps the arc of $([u,u'],[n,n'])$ onto the arc of the same
dual pair, the vertex $p_q$ of the discriminant onto the zero-cell $C^\Upsilon([u,u'],P_q)$, the negative vertices of $\Gamma$ (edge,
triangle) onto the vertices of $\Gamma^\Upsilon$ of positive type, and the positive vertices (triangle, edge) onto the
vertices of negative type. At a point of an open facet $\operatorname{int}F_u$, where $\Lambda=u^\perp\cap N$ and
$\Lambda^\Upsilon=M/\Z u$ at the image point, $\kappa_{\mathfrak h}(x)=\langle\,\cdot\,,x\rangle$.
\item On the arc $A=C^\Upsilon([u,u'],[n,n'])$: $u'-u$ is a flat section of $\Lambda^\Upsilon$ near $A$;
$d^*_A=\langle\,\cdot\,,n'-n\rangle$ is a flat section of $(\Lambda^\Upsilon)^\vee$ near $A$, and $d^*_A=\kappa_{\mathfrak h}(n'-n)$;
\[
\operatorname{Inv}^*(A)=\{X\in(\Lambda^\Upsilon)^\vee:X(u'-u)=0\},
\]
of rank two, with $d^*_A$ primitive in it; and the sublattice of $\Lambda^\Upsilon$ fixed by the monodromy around $A$ is
$\{n,n'\}^\perp\cap M=\ker d^*_A$, of rank two.
\item At the vertices of $\Gamma^\Upsilon$:
\begin{itemize}
\item at a node $C^\Upsilon([u,u'],P)$, $P=abcd$: $\operatorname{Inv}^*=\{X:X(u'-u)=0\}=\operatorname{Inv}^*(A)$ for each
of the four arcs $A$ at it, with basis $\langle\,\cdot\,,b-a\rangle$, $\langle\,\cdot\,,c-b\rangle$;
\item at a vertex $C^\Upsilon([u,u'],\{n_a,n_b,n_c\})$ of positive type: $\operatorname{Inv}^*=\{X:X(u'-u)=0\}$, of rank
two, containing $d^*_A$ for the three arcs $A$ at it;
\item at a vertex $C^\Upsilon(\{u_1,u_2,u_3\},[n,n'])$ of negative type: $\operatorname{Inv}^*=\Z\langle\,\cdot\,,n'-n\rangle$,
the common value of $d^*_A$ on the three arcs at it, up to sign.
\end{itemize}
In each case $\kappa_{\mathfrak h}$ maps the stalk of $\iota_*\Lambda$ at the corresponding point of $\Gamma$ onto this group.
\item For a $1$-cycle $\sum_Ac_A\,d^*_A[A]$ of $C_*(\Gamma^\Upsilon;\mathcal D^*)$, with $d^*_A=\langle\,\cdot\,,n'_A-n_A\rangle$,
define its canonical lift $\sum_Ac_A(e_{n'_A}-e_{n_A})[A]\in C_1(\Gamma^\Upsilon;\Z^R)$ and its node coefficients $\rho^*_q$
by the rule of Definition~\ref{def:nodecoef} at the zero-cell $C^\Upsilon([u,u'],P_q)$. Equivalently, with the coefficients
oriented away from the node and written $c_{ab}\langle\,\cdot\,,b-a\rangle$ and $c_{bc}\langle\,\cdot\,,c-b\rangle$ on the
arcs over $ab$ and $bc$, $\rho^*_q=c_{ab}-c_{bc}$.
\end{enumerate}
\end{proposition}
\begin{proof}
(1) The first statement combines Proposition~\ref{prop:F}(3) for $g_C$, Proposition~\ref{prop:L1}(4) for $\mathcal L$ and
the definition of $C^\Upsilon({\mathfrak p})$; the vertex types are those of Proposition~\ref{prop:L2}(3), and $\Upsilon$ carries the
monodromies of $\Lambda_Y$ to those of $\Lambda^\Upsilon$. The map $g_C$ preserves $\operatorname{int}F_u$, where
$\kappa_C$ is the identity of $u^\perp\cap N$, and $\mathcal L(\operatorname{int}F_u)=\widetilde R^*_u$, where
$\Lambda_Y=M/\Z u$ and the flat pairing is $\langle\bar\zeta,x\rangle_{\mathcal L}=\langle\zeta,x\rangle$ (proof of Proposition~\ref{prop:L3}).

(2) Transport by $\Upsilon$ does not change monodromies, so it suffices to argue on $B_Y$ near the arc $C_Y([u,u'],[n,n'])$.
In the chart of $C_Y(u,n)$, $\Lambda_Y=n^\perp\cap M$, which contains $u'-u$ because $\langle u,n\rangle=\langle
u',n\rangle=-1$. By Proposition~\ref{prop:L2}(4) the monodromy is $T(\zeta)=\zeta+\langle\zeta,n'-n\rangle(u'-u)$. So
$T(u'-u)=u'-u$, as $\langle u'-u,n'-n\rangle=0$; the fixed sublattice is $\{\zeta\in n^\perp\cap M:\langle
\zeta,n'-n\rangle=0\}=\{n,n'\}^\perp\cap M$, of rank two because $n,n'$ are linearly independent; and a covector $X$ is
invariant under the dual action if and only if $X(T\zeta)=X(\zeta)$ for all $\zeta$, that is, $\langle\zeta,n'-n\rangle X(u'-u)=0$
for all $\zeta\in n^\perp\cap M$, that is, $X(u'-u)=0$, since $\langle\,\cdot\,,n'-n\rangle$ does not vanish on $n^\perp\cap M$.
The covector $\langle\,\cdot\,,n'-n\rangle$ vanishes on $u'-u$, hence is flat, and it is primitive in $N/\Z n$
(proof of Proposition~\ref{prop:L2}(3)). By (1), $\kappa_{\mathfrak h}(n'-n)=\langle\,\cdot\,,n'-n\rangle$ on the side of $F_u$, and both
are flat sections near the arc.

(3) At a node or at a vertex of positive type, all arcs share the vector $u'-u$ (Proposition~\ref{prop:L2}(3)) and have
monodromies $\zeta\mapsto\zeta+\langle\zeta,t_A\rangle(u'-u)$ with $t_A$ the edge vectors of $P$, respectively of
the triangle; as in (2) the invariant covectors are those vanishing on $u'-u$. At a node, in the chart of $C_Y(u,a)$,
$(\Lambda_Y)^\vee=N/\Z a$ and the covectors vanishing on $u'-u$ form the image of $(u'-u)^\perp\cap N$. The plane of $G$ is
at lattice distance one from the origin and the triangle $abd$ is unimodular in it, so $\{a,b-a,d-a\}$ is a basis of
$(u'-u)^\perp\cap N$, and the classes of $b-a$ and $d-a=c-b$ form a basis of the invariants. At a vertex of negative type
the three arcs share the covector $\langle\,\cdot\,,n'-n\rangle$ and have monodromies $\zeta\mapsto\zeta+\langle
\zeta,n'-n\rangle(u_j-u_i)$; a covector is invariant if and only if it vanishes on the $u_j-u_i$, which span a saturated
rank-two sublattice of $n^\perp\cap M$ (proof of Proposition~\ref{prop:L2}(3)); its annihilator in $N/\Z n$ is saturated
of rank one and contains the primitive class of $n'-n$. The last statement holds because $\kappa_{\mathfrak h}$ is an isomorphism of local
systems and $\mathfrak h$ maps small balls about a point of $\Gamma$ onto neighbourhoods of its image.

(4) The equivalence of the two descriptions is the computation of Section~\ref{ssec:cycles-bdry} for
Definition~\ref{def:nodecoef}, which uses only the labels $n_A,n'_A$ of the arcs at the node.
\end{proof}

\begin{proposition}[transfer of relative tropical $2$-cycles]\label{prop:transfer}
Let $\mathfrak K$ be a triangulation of $B'_\sigma$ in which $\Gamma$ is full, and $\mathfrak K_1$ as in
Definition~\ref{def:transfer}. Then $\mathfrak h(\mathfrak K_1)$ is a triangulation of $B_Y$ in which $\Gamma^\Upsilon$ is full, and
\[
\mathfrak h_*\colon Z(\mathfrak K)\longrightarrow Z\bigl(\mathfrak h(\mathfrak K_1);F^*,\mathcal D^*\bigr)
\]
is an injective homomorphism; if $\mathfrak h$ is linear on the simplices of $\mathfrak K$ and $\mathfrak K_1=\mathfrak K$, it is an
isomorphism. If the boundary network of $z$ is
$\sum_Ec_E\,d_E[E]$, the boundary network of $\mathfrak h_*z$ is $\sum_Ec_E\sum_{E'\subseteq E}d^*_{\mathfrak h(A(E))}[\mathfrak h(E')]$, with the same
integers $c_E$, and
\[
\rho^*(\mathfrak h_*z)=\rho(z).
\]
\end{proposition}
\begin{proof}
By Proposition~\ref{prop:dict}(1), $\mathfrak h$ is a piecewise linear homeomorphism with $\mathfrak h(\Gamma)=\Gamma^\Upsilon$ mapping arcs
onto arcs, and by Proposition~\ref{prop:dict}(2), $\kappa_{\mathfrak h}(d_E)=\kappa_{\mathfrak h}(n'-n)=d^*_{\mathfrak h(A)}$ for a leg $E$ with lower support
$\{n,n'\}$ on the arc $A$ of $([u,u'],[n,n'])$ (Proposition~\ref{prop:L2}(2)). So $(\mathfrak h,\kappa_{\mathfrak h})$ satisfies the hypotheses of
Lemma~\ref{lem:transport}, which gives all but the last statement. The node coefficients of $z$ are computed from the
integers $c_E$ on the legs at $p_q$ and the lower supports of these legs; those of $\mathfrak h_*z$ from the same integers on the arcs
at $\mathfrak h(p_q)=C^\Upsilon([u,u'],P_q)$ and the same labels (Proposition~\ref{prop:dict}(1),(4)). Subdivision replaces each leg
by edges with the same coefficient, which does not change the boundary of the canonical lift at the node.
\end{proof}

Different choices of $\mathfrak K^0_1$ and $\mathfrak K_1$ give transfers with the same boundary networks, up to
subdivision, and the same node coefficients; no orientation of $B'_\sigma$ or of $B_Y$ enters the definition. By Proposition~\ref{prop:dict}(3) the transferred coefficient system has, at
the vertices of $\Gamma^\Upsilon$, the stalks of $\iota_*\Lambda$ at the corresponding vertices of $\Gamma$, with the types
exchanged: the rank-one stalks, which sit at the positive vertices of $B'_\sigma$, sit on $B_Y$ at the vertices of negative
type, and the rank-two stalks of the negative vertices of $B'_\sigma$ at the vertices of positive type.

\begin{remark}[tropical $2$-cycles and the lift of Casta\~no-Bernard and Matessi]\label{rem:CBMfibre}
Let $(S,j,v)$ be a tropical $2$-cycle in the sense of \cite[Def.~7.2]{CBM}, triangulated, so that it defines a relative
tropical $2$-cycle $z$ with coefficient $v$ on the triangles of its smooth part (Section~\ref{ssec:cbm}). The coefficients
of $\mathfrak h_*z$ are the covectors $\kappa_{\mathfrak h}(v)$, and $\ker\kappa_{\mathfrak h}(v)\subseteq\Lambda^\Upsilon$ is a rank-two subbundle along the image of
the smooth part. It is the bundle $\ker v\subseteq T^*B$ of \cite[Def.~7.2]{CBM}, written in the affine structure of $B_Y$.
In \cite[\S7]{CBM} the lift $\widetilde S$ of $S$ lies in the torus fibration over $B'_\sigma$ with fibres
$T^*_bB/\Lambda^\vee_b$, with fibre $(\ker v)_b/(\ker v)_b\cap\Lambda^\vee_b$ over a smooth point; by
Proposition~\ref{prop:L3} this is the torus fibration over $B_Y$ with fibres $\Lambda_Y\otimes U(1)$, and the fibre of the
lift over the image point is the two-torus $\ker\kappa_{\mathfrak h}(v)\otimes U(1)$. This is the fibre of the chain of
Theorem~\ref{thm:lift} over the regular part.
\end{remark}

\begin{remark}[why a transfer map is needed]\label{rem:mismatch}
A relative tropical $2$-cycle on $B'_\sigma$ is not read directly as an object on $B_Y$, for two reasons. First, the
polyhedral decompositions need not correspond. An inclusion-reversing bijection $c\mapsto\check
c$ from the cells of $\mathscr P_X$ onto those of $\mathscr P_Y$ with $\tau_Y(\check c)=\tau(c)$ and $\omega_Y(\check
c)=\omega(c)$ would map the legs of the arc $C_X(\tau,\omega)$, the chains $e\subsetneq f$ with $\omega(e)=\omega$ and
$\tau(f)=\tau$, bijectively onto the legs of $C_Y(\tau,\omega)$, the chains $f'\subsetneq e'$ with $\tau_Y(f')=\tau$ and
$\omega_Y(e')=\omega$, so corresponding arcs would consist of the same number of legs. These numbers are determined by
mixed subdivisions of different faces with different heights, of $F_\tau$ with $h$ on $\tau^\vee$ for $C_X(\tau,\omega)$
and of $F^*_\omega$ with $\check h$ on $\omega^\vee$ for $C_Y(\tau,\omega)$; the dual-pair complex needs no such
coincidence. Second, the transfer of a tropical $2$-cycle is not a tropical $2$-cycle of $B_Y$ in the sense of \cite[Def.~7.2]{CBM}: its coefficients are covectors; its boundary
vertices, which lie at negative vertices of $B'_\sigma$, lie at vertices of positive type of $B_Y$; and at a node met by its
boundary, which is of positive type on $B_Y$ by Proposition~\ref{prop:L2}(3), its boundary continues along one line of the
discriminant, whereas \cite[Def.~7.2(f)]{CBM} requires the boundary to turn from one line to the other at a positive node
\cite[Fig.~14(3)]{CBM}. The transfer $\mathfrak h_*$ records the data from which the lift is built, in the affine structure of $B_Y$,
without asking them to be a tropical object there.
\end{remark}
